\ifdefined\gkver\else\def\gkver{student}\fi
\documentclass[\gkver]{gaokaozhenti}
\begin{document}

\chapter{2009年全国理综Ⅰ}
%% paperType: 理综物理部分

\begin{enumerate}
\item
%% number: 14
%% paperName: 2009年全国理综Ⅰ第 14 题 6 分
%% typeId: 1
%% type: 单选题
%% chapter: 气体性质
%% point: 气体的状态参量
%% method: 无
%% score: 6
%% degree: 450
%% duplicateId: 0
%% body:
下列说法正确的是\xzanswer{A}

\fourchoices[answer=A]
{气体对器壁的压强就是大量气体分子作用在器壁单位面积上的平均作用力}
{气体对器壁的压强就是大量气体分子单位时间作用在器壁上的平均冲量}
{气体分子热运动的平均动能减小，气体的压强一定减小}
{单位体积的气体分子数增加，气体的压强一定增大}

%% answer: A
%% memo:
\memoanswer{
本题考查气体部分的知识。根据压强的定义 A 正确，B 错。气体分子热运动的平均动能减小，说明温度降低，但不能说明压强也一定减小，C 错。单位体积的气体分子增加，但温度降低有可能气体的压强减小，D 错。
}

\item
%% number: 15
%% paperName: 2009年全国理综Ⅰ第 15 题 6 分
%% typeId: 1
%% type: 单选题
%% chapter: 光学
%% point: 光的折射
%% method: 无
%% score: 6
%% degree: 650
%% duplicateId: 0
%% body:
某物体左右两侧各有一竖直放置的平面镜，两平面镜相互平行，物体距左镜 $4\Um$，右镜 $8\Um$，如图所示，物体在左镜所成的像中从右向左数的第三个像与物体的距离是\xzanswer{B}

\onepicture[w=5.0cm]{figs/15.png}

\fourchoices[answer=B]
{$24\Um$}
{$32\Um$}
{$40\Um$}
{$48\Um$}

%% answer: B
%% memo:
\memoanswer{
本题考查平面镜成像。从右向左在左镜中的第一个像是物体的像，距离物体 $8\Um$，第二个像是物体在右镜所成像的像，第三个像是第一个像在右镜中的像在左镜中的像，距离物体 $32\Um$。
}

\item
%% number: 16
%% paperName: 2009年全国理综Ⅰ第 16 题 6 分
%% typeId: 1
%% type: 单选题
%% chapter: 原子和原子核
%% point: 玻尔模型
%% method: 无
%% score: 6
%% degree: 650
%% duplicateId: 0
%% body:
氦氖激光器能产生三种波长的激光，其中两种波长分别为 $\lambda_{1}=0.6328\Uum$，$\lambda_{2}=3.39\Uum$。已知波长为 $\lambda_{1}$ 的激光是氖原子在能级间隔为 $\Delta E_{1}=1.96\UeV$ 的两个能级之间跃迁产生的。用 $\Delta E_{2}$ 表示产生波长为 $\lambda_{2}$ 的激光所对应的跃迁的能级间隔，则 $\Delta E_{2}$ 的近似值为\xzanswer{D}

\fourchoices[answer=D]
{$10.50\UeV$}
{$0.98\UeV$}
{$0.53\UeV$}
{$0.36\UeV$}

%% answer: D
%% memo:
\memoanswer{
本题考查玻尔的原子跃迁理论。根据 $\Delta E=h\nu$，$\nu=\frac{c}{\lambda}$，可知当 $\Delta E_{1}=1.96\UeV$，$\lambda_{1}=0.6328\Uum$ 时，当 $\lambda_{2}=3.39\Uum$ 时，联立可知 $\Delta E_{2}=0.36\UeV$。
}

\item
%% number: 17
%% paperName: 2009年全国理综Ⅰ第 17 题 6 分
%% typeId: 1
%% type: 单选题
%% chapter: 磁场
%% point: 安培力
%% method: 等效替代
%% score: 6
%% degree: 650
%% duplicateId: 0
%% body:
如图，一段导线 abcd 位于磁感应强度大小为 $B$ 的匀强磁场中，且与磁场方向（垂直于纸面向里）垂直。线段 ab、bc 和 cd 的长度均为 $L$，且 $\angle abc=\angle bcd=135^\circ$。流经导线的电流为 $I$，方向如图中箭头所示。导线段 abcd 所受到的磁场的作用力的合力\xzanswer{A}

\onepicture[w=6.0cm]{figs/17.png}

\fourchoices[answer=A]
{方向沿纸面向上，大小为 $(\sqrt{2}+1)IBL$}
{方向沿纸面向上，大小为 $(\sqrt{2}-1)IBL$}
{方向沿纸面向下，大小为 $(\sqrt{2}+1)IBL$}
{方向沿纸面向下，大小为 $(\sqrt{2}-1)IBL$}

%% answer: A
%% memo:
\memoanswer{
本题考查安培力的大小与方向的判断。该导线可以用 a 和 d 之间的直导线长为 $(\sqrt{2}+1)L$ 来等效代替，根据 $F=BIl$，可知大小为 $(\sqrt{2}+1)IBL$，方向根据左手定则判断，A 正确。
}

\item
%% number: 18
%% paperName: 2009年全国理综Ⅰ第 18 题 6 分
%% typeId: 2
%% type: 多选题
%% chapter: 电场
%% point: 电场线
%% method: 无
%% score: 6
%% degree: 450
%% duplicateId: 0
%% body:
如图所示，一电场的电场线关于 $y$ 轴（沿竖直方向）对称，O、M、N 是 $y$ 轴上的三个点，且 $OM=MN$。则\xzanswer{AD}

\onepicture[w=6.0cm]{figs/18.png}

\fourchoices[answer=AD]
{M 点的电势比 P 点的电势高}
{将负电荷由 O 点移动到 P 点，电场力做正功}
{M、N 两点间的电势差大于 O、M 两点间的电势差}
{在 O 点静止释放一带正电粒子，该粒子将沿 $y$ 轴做直线运动}

%% answer: AD
%% memo:
\memoanswer{
本题考查电场、电势、等势线以及带电粒子在电场中的运动。由图和几何关系可知 M 和 P 两点不处在同一等势线上，而且有 $\varphi_{M}>\varphi_{P}$，A 对；将负电荷由 O 点移到 P 要克服电场力做功，即电场力做负功，B 错；根据 $U=Ed$，O 到 M 的平均电场强度大于 M 到 N 的平均电场强度，所以有 $U_{OM}>U_{MN}$，C 错；从 O 点释放正电荷后，电场力做正功，该粒子将沿 $y$ 轴做加速直线运动，D 对。
}

\item
%% number: 19
%% paperName: 2009年全国理综Ⅰ第 19 题 6 分
%% typeId: 1
%% type: 单选题
%% chapter: 曲线运动
%% point: 天体运动
%% method: 比例法
%% score: 6
%% degree: 450
%% duplicateId: 0
%% body:
天文学家新发现了太阳系外的一颗行星。这颗行星的体积是地球的 4.7 倍，质量是地球的 25 倍，已知某一近地卫星绕地球运动的周期约为 $1.4\Uh$，引力常量 $G=6.67\times10^{-11}\UN\cdot\Umq/\Ukg^{2}$，由此估算该行星的平均密度约为\xzanswer{D}

\fourchoices[answer=D]
{$1.8\times10^{3}\Ukgmc$}
{$5.6\times10^{3}\Ukgmc$}
{$1.1\times10^{4}\Ukgmc$}
{$2.9\times10^{4}\Ukgmc$}

%% answer: D
%% memo:
\memoanswer{
本题考查天体运动的知识。首先根据近地卫星绕地球运动的向心力由万有引力提供 $G\frac{Mm}{R^{2}}=m\frac{4\pi^{2}R}{T^{2}}$，可求出地球的质量，然后根据 $\rho=\frac{3M}{4\pi R^{3}}$，可得该行星的密度约为 $2.9\times10^{4}\Ukgmc$。
}

\item
%% number: 20
%% paperName: 2009年全国理综Ⅰ第 20 题 6 分
%% typeId: 2
%% type: 多选题
%% chapter: 机械振动
%% point: 振动图象
%% method: 无
%% score: 6
%% degree: 450
%% duplicateId: 0
%% body:
一列简谐横波在某一时刻的波形图如图 \ref{2009全国理综Ⅰ20a} 所示，图中 P、Q 两质点的横坐标分别为 $x=1.5\Um$ 和 $x=4.5\Um$。P 点的振动图象如图 \ref{2009全国理综Ⅰ20b} 所示。

\twopicture[labelA=2009全国理综Ⅰ20a, labelB=2009全国理综Ⅰ20b, numA=1, numB=2, w=6.2cm]
{figs/20a.png}
{figs/20b.png}

在下列四幅图中，Q 点的振动图象可能是\xzanswer{BC}

\fourchoices[answer=BC, ispicture=true, h=2.2cm]
{figs/20c.png}
{figs/20d.png}
{figs/20e.png}
{figs/20f.png}

%% answer: BC
%% memo:
\memoanswer{
本题考查波的传播。该波的波长为 $4\Um$，P、Q 两点间的距离为 $3\Um$。当波沿 $x$ 轴正方向传播时，P 在平衡位置向上振动，而 Q 点此时应处于波峰，B 正确；当沿 $x$ 轴负方向传播时，P 点处于向上振动时 Q 点应处于波谷，C 对。
}

\item
%% number: 21
%% paperName: 2009年全国理综Ⅰ第 21 题 6 分
%% typeId: 2
%% type: 多选题
%% chapter: 运动和力
%% point: 动量守恒定律
%% method: 无
%% score: 6
%% degree: 450
%% duplicateId: 0
%% body:
质量为 $M$ 的物块以速度 $v$ 运动，与质量为 $m$ 的静止物块发生正碰，碰撞后两者的动量正好相等，两者质量之比 $M/m$ 可能为\xzanswer{AB}

\fourchoices[answer=AB]
{2}
{3}
{4}
{5}

%% answer: AB
%% memo:
\memoanswer{
本题考查动量守恒。设碰撞后两者的动量都为 $P$，根据动量守恒和能量守恒得，总动量为 $2P$，根据 $P^{2}=2mE_{K}$，以及能量的关系得 $\frac{(2P)^{2}}{2M}\geq\frac{P^{2}}{2m}+\frac{P^{2}}{2M}$，得 $\frac{M}{m}\leq3$，所以 A、B 正确。
}

\item
%% number: 22
%% paperName: 2009年全国理综Ⅰ第 22 题 8 分
%% typeId: 4
%% type: 实验题
%% chapter: 电路
%% point: 故障分析
%% method: 无
%% score: 8
%% degree: 450
%% duplicateId: 0
%% body:
如图所示的电路，1、2、3、4、5、6 为连接点的标号。在开关闭合后，发现小灯泡不亮。现用多用电表检查电路故障，需要检测的有：电源、开关、小灯泡、3 根导线以及电路中的各连接点。

\begin{enumerate}
	\item
	为了检测小灯泡以及 3 根导线，在连接点 1、2 已接好的情况下，应当选用多用电表的 \tkanswer[1.2cm]{} 档，在连接点 1、2 同时断开的情况下，应当选用多用电表的 \tkanswer[1.2cm]{} 档。
	\item
	在开关闭合情况下，若测得 5、6 两点间的电压接近电源的电动势，则表明 \tkanswer[4cm]{} 可能有故障。
	\item
	将小灯泡拆离电路，写出用多用电表检测小灯泡是否有故障的具体步骤。
\end{enumerate}
\onepicture[w=5.0cm, align=right]{figs/22.png}

%% answer:
\jdanswer{
\begin{enumerate}
	\item
	电压；欧姆
	\item
	开关或连接点 5、6
	\item
	①调到欧姆档；②将红、黑表笔相接，检查欧姆档能否正常工作；③测量小灯泡的电阻，如电阻无穷大，表明小灯泡有故障。
\end{enumerate}
}

%% memo:
\memoanswer{
（1）在 1、2 两点接好的情况下，应当选用多用电表的电压档；在 1、2 同时断开的情况下，应选用欧姆档。（2）表明 5、6 两点可能有故障。（3）①调到欧姆档；②将红、黑表笔相接，检查欧姆档能否正常工作；③测量小灯泡的电阻，如电阻无穷大，表明小灯泡有故障。
}

\item
%% number: 23
%% paperName: 2009年全国理综Ⅰ第 23 题 10 分
%% typeId: 4
%% type: 实验题
%% chapter: 运动和力
%% point: 验证牛顿第二定律
%% method: 无
%% score: 10
%% degree: 450
%% duplicateId: 0
%% body:
某同学为了探究物体在斜面上运动时摩擦力与斜面倾角的关系，设计实验装置如图。长直平板一端放在水平桌面上，另一端架在一物块上。在平板上标出 A、B 两点，B 点处放置一光电门，用光电计时器记录滑块通过光电门时挡光的时间，实验步骤如下：

\begin{steps}
	\item
	用游标卡尺测量滑块的挡光长度 $d$，用天平测量滑块的质量 $m$；
	\item
	用直尺测量 A、B 之间的距离 $s$，A 点到水平桌面的垂直距离 $h_{1}$，B 点到水平桌面的垂直距离 $h_{2}$；
	\item
	将滑块从 A 点静止释放，由光电计时器读出滑块的挡光时间 $t$；
	\item
	重复步骤③数次，并求挡光时间的平均值 $\overline{t}$；
	\item
	利用所测数据求出摩擦力 $f$ 和斜面倾角的余弦值 $\cos\alpha$；
	\item
	多次改变斜面的倾角，重复实验步骤②③④⑤做出 $f-\cos\alpha$ 关系曲线。
\end{steps}

\begin{enumerate}
	\item
	用测量的物理量完成下列各式（重力加速度为 $g$）

	①斜面倾角的余弦 $\cos\alpha=$ \tkanswer[2cm]{}；

	②滑块通过光电门时的速度 $v=$ \tkanswer[2cm]{}；

	③滑块运动时的加速度 $a=$ \tkanswer[2cm]{}；

	④滑块运动时所受到的摩擦阻力 $f=$ \tkanswer[2cm]{}。
	\item
	测量滑块挡光长度的游标卡尺读数如图所示，读得 $d=$ \tkanswer[2cm]{}。
\end{enumerate}
\twopicture[labelA=2009全国理综Ⅰ23a, labelB=2009全国理综Ⅰ23b, w=6.2cm, align=right]
{figs/23a.png}
{figs/23b.png}

%% answer:
\jdanswer{
\begin{enumerate}
	\item
	①$\frac{1}{s}\sqrt{s^{2}-(h_{1}-h_{2})^{2}}$；②$\frac{d}{\overline{t}}$；③$\frac{d^{2}}{2s\overline{t}^{2}}$；④$mg\frac{h_{1}-h_{2}}{s}-m\frac{d^{2}}{2s\overline{t}^{2}}$。
	\item
	$3.62\Ucm$。
\end{enumerate}
}

%% memo:
\memoanswer{
（1）物块在斜面上做初速度为零的匀加速直线运动，受重力、支持力、滑动摩擦力。①根据三角形关系可得到 $\cos\alpha=\frac{\sqrt{s^{2}-(h_{1}-h_{2})^{2}}}{s}$；②根据 $v=\frac{x}{t}=\frac{d}{\overline{t}}$；③根据运动学公式 $x=\frac{v^{2}}{2a}$，有 $s=\frac{v^{2}}{2a}$，即有 $a=\frac{d^{2}}{2s\overline{t}^{2}}$；④根据牛顿第二定律 $mg\sin\theta-f=ma$，则有 $f=mg\frac{h_{1}-h_{2}}{s}-m\frac{d^{2}}{2s\overline{t}^{2}}$。（2）在游标卡尺中，主尺上是 $3.6\Ucm$，在游标尺上恰好是第 2 条刻度线与主尺对齐，再考虑到卡尺是 10 分度，所以读数为 $3.6\Ucm+0.2\times1\Umm=3.62\Ucm$。
}

\item
%% number: 24
%% paperName: 2009年全国理综Ⅰ第 24 题 15 分
%% typeId: 6
%% type: 计算题
%% chapter: 电路
%% point: 电阻定律
%% method: 无
%% score: 15
%% degree: 450
%% duplicateId: 0
%% body:
材料的电阻率 $\rho$ 随温度变化的规律为 $\rho=\rho_{0}(1+\alpha t)$，其中 $\alpha$ 称为电阻温度系数，$\rho_{0}$ 是材料在 $t=0\UdC$ 时的电阻率。在一定温度范围内 $\alpha$ 是与温度无关的常量。金属的电阻一般随温度的增加而增加，具有正温度系数；而某些非金属如碳等则相反，具有负温度系数。利用具有正负温度系数的两种材料的互补特性，可制成阻值在一定温度范围内不随温度变化的电阻。已知在 $0\UdC$ 时，铜的电阻率为 $1.7\times10^{-8}\UO\cdot\Um$，碳的电阻率为 $3.5\times10^{-5}\UO\cdot\Um$；在 $0\UdC$ 附近，铜的电阻温度系数为 $3.9\times10^{-3}\UdC^{-1}$，碳的电阻温度系数为 $-5\times10^{-4}\UdC^{-1}$。将横截面积相同的碳棒与铜棒串接成长 $1.0\Um$ 的导体，要求其电阻在 $0\UdC$ 附近不随温度变化，求所需碳棒的长度（忽略碳棒和铜棒的尺寸随温度的变化）。

%% answer:
\jdanswer{$3.8\times10^{-3}\Um$}

%% memo:
\memoanswer{
设所需碳棒的长度为 $L_{1}$，电阻率为 $\rho_{1}$，电阻温度系数为 $\alpha_{1}$；铜棒的长度为 $L_{2}$，电阻率为 $\rho_{2}$，电阻温度系数为 $\alpha_{2}$。根据题意有 $\rho_{1}=\rho_{10}(1+\alpha_{1}t)$ ①，$\rho_{2}=\rho_{20}(1+\alpha_{2}t)$ ②，式中 $\rho_{10}$、$\rho_{20}$ 分别为碳和铜在 $0\UdC$ 时的电阻率。

设碳棒的电阻为 $R_{1}$，铜棒的电阻为 $R_{2}$，有 $R_{1}=\rho_{1}\frac{L_{1}}{S}$ ③，$R_{2}=\rho_{2}\frac{L_{2}}{S}$ ④，式中 $S$ 为碳棒与铜棒的横截面积。

碳棒和铜棒连接成的导体的总电阻和总长度分别为 $R=R_{1}+R_{2}$ ⑤，$L_{0}=L_{1}+L_{2}$ ⑥，式中 $L_{0}=1.0\Um$。

联立以上各式得 $R=\rho_{10}\frac{L_{1}}{S}+\rho_{20}\frac{L_{2}}{S}+\frac{\rho_{10}\alpha_{1}L_{1}+\rho_{20}\alpha_{2}L_{2}}{S}t$ ⑦。

要使电阻 $R$ 不随温度 $t$ 变化，⑦式中 $t$ 的系数必须为零，即 $\rho_{10}\alpha_{1}L_{1}+\rho_{20}\alpha_{2}L_{2}=0$ ⑧。

联立⑥⑧得 $L_{1}=\frac{\rho_{20}\alpha_{2}}{\rho_{20}\alpha_{2}-\rho_{10}\alpha_{1}}L_{0}$ ⑨，代入数据解得 $L_{1}=3.8\times10^{-3}\Um$ ⑩。
}

\item
%% number: 25
%% paperName: 2009年全国理综Ⅰ第 25 题 18 分
%% typeId: 6
%% type: 计算题
%% chapter: 运动和力
%% point: 动量和能量
%% method: 无
%% score: 18
%% degree: 450
%% duplicateId: 0
%% body:
如图所示，倾角为 $\theta$ 的斜面上静止放置三个质量均为 $m$ 的木箱，相邻两木箱的距离均为 $l$。工人用沿斜面的力推最下面的木箱使之上滑，逐一与其它木箱碰撞。每次碰撞后木箱都粘在一起运动。整个过程中工人的推力不变，最后恰好能推着三个木箱匀速上滑。已知木箱与斜面间的动摩擦因数为 $\mu$，重力加速度为 $g$。设碰撞时间极短，求：

\begin{enumerate}
	\item
	工人的推力；
	\item
	三个木箱匀速运动的速度；
	\item
	在第一次碰撞后损失的机械能。
\end{enumerate}
\onepicture[w=6.0cm, align=right]{figs/25.png}

%% answer:
\jdanswer{
\begin{enumerate}
	\item
	$3mg(\sin\theta+\mu\cos\theta)$
	\item
	$\frac{2}{3}\sqrt{2gl(\sin\theta+\mu\cos\theta)}$
	\item
	$\Delta E=mgl(\sin\theta+\mu\cos\theta)$
\end{enumerate}
}

%% memo:
\memoanswer{
（1）当匀速时，把三个物体看作一个整体受重力、推力 $F$、摩擦力 $f$ 和支持力，根据平衡的知识有 $F=3mg\sin\theta+3\mu mg\cos\theta=3mg(\sin\theta+\mu\cos\theta)$。

（2）第一个木箱与第二个木箱碰撞之前的速度为 $v_{1}$，加速度 $a_{1}=\frac{F-mg\sin\theta-\mu mg\cos\theta}{m}=2g(\sin\theta+\mu\cos\theta)$。根据运动学公式或动能定理有 $v_{1}=2\sqrt{gl(\sin\theta+\mu\cos\theta)}$，碰撞后的速度为 $v_{2}$，根据动量守恒有 $mv_{1}=2mv_{2}$，即碰撞后的速度为 $v_{2}=\sqrt{gl(\sin\theta+\mu\cos\theta)}$。然后一起去碰撞第三个木箱，设碰撞前的速度为 $v_{3}$，从 $v_{2}$ 到 $v_{3}$ 的加速度为 $a_{2}=\frac{F-2mg\sin\theta-2\mu mg\cos\theta}{2m}=\frac{g(\sin\theta+\mu\cos\theta)}{2}$，根据运动学公式有 $v_{3}^{2}-v_{2}^{2}=2a_{2}l$，得 $v_{3}=\sqrt{2gl(\sin\theta+\mu\cos\theta)}$，跟第三个木箱碰撞根据动量守恒有 $2mv_{3}=3mv_{4}$，得 $v_{4}=\frac{2}{3}\sqrt{2gl(\sin\theta+\mu\cos\theta)}$，就是匀速的速度。

（3）设第一次碰撞中的能量损失为 $\Delta E$，根据能量守恒有 $\frac{1}{2}mv_{1}^{2}=\Delta E+\frac{1}{2}\times2mv_{2}^{2}$，代入数据得 $\Delta E=mgl(\sin\theta+\mu\cos\theta)$。
}

\item
%% number: 26
%% paperName: 2009年全国理综Ⅰ第 26 题 21 分
%% typeId: 6
%% type: 计算题
%% chapter: 磁场
%% point: 带电粒子在磁场中的运动
%% method: 无
%% score: 21
%% degree: 250
%% duplicateId: 0
%% body:
如图，在 $x$ 轴下方有匀强磁场，磁感应强度大小为 $B$，方向垂直于 $xy$ 平面向外，P 是 $y$ 轴上距原点为 $h$ 的一点，$N_{0}$ 为 $x$ 轴上距原点为 $a$ 的一点。A 是一块平行于 $x$ 轴的挡板，与 $x$ 轴的距离为 $\frac{h}{2}$，A 的中点在 $y$ 轴上，长度略小于 $\frac{a}{2}$。带电粒子与挡板碰撞前后 $x$ 方向上的分速度不变，$y$ 方向上的分速度反向，大小不变。质量为 $m$、电荷量为 $q$（$q>0$）的粒子从 P 点瞄准 $N_{0}$ 点入射，最后又通过 P 点。不计重力。求粒子入射速度的所有可能值。

\onepicture[w=6.0cm, align=right]{figs/26.png}

%% answer:
\jdanswer{
\begin{enumerate}
	\item
	$v_{0}=\frac{qBa\sqrt{a^{2}+h^{2}}}{mh}$，$n=0$
	\item
	$v_{1}=\frac{3qBa\sqrt{a^{2}+h^{2}}}{4mh}$，$n=1$
	\item
	$v_{2}=\frac{2qBa\sqrt{a^{2}+h^{2}}}{3mh}$，$n=2$
\end{enumerate}
}

%% memo:
\memoanswer{
设粒子的入射速度为 $v$，第一次射出磁场的点为 $N_{0}'$，与板碰撞后再次进入磁场的点为 $N_{1}$。粒子在磁场中运动的半径为 $R$，有 $R=\frac{mv}{qB}$，粒子速率不变，每次射入磁场与射出磁场的点间距 $x_{1}$ 不变，$x_{1}=N_{0}'N_{0}=2R\sin\theta$；粒子射出磁场与下一次射入磁场的点间距 $x_{2}$ 不变，$x_{2}=a$。设最终离开磁场时与挡板碰撞 $n$ 次，若粒子能绕回到 P 点，由对称性可知，最后一次出射点的坐标应为 $-a$，即 $(n+1)x_{1}-nx_{2}=2a$，联立解得 $x_{1}=\frac{n+2}{n+1}a$。若粒子与挡板发生碰撞，有 $x_{1}-x_{2}>\frac{a}{4}$，所以 $n<3$，$v=\frac{qB}{2m\sin\theta}\cdot\frac{n+2}{n+1}a$，式中 $\sin\theta=\frac{h}{\sqrt{a^{2}+h^{2}}}$，则：
$v_{0}=\frac{qBa\sqrt{a^{2}+h^{2}}}{mh}$，$n=0$；
$v_{1}=\frac{3qBa\sqrt{a^{2}+h^{2}}}{4mh}$，$n=1$；
$v_{2}=\frac{2qBa\sqrt{a^{2}+h^{2}}}{3mh}$，$n=2$。

\onepicture[w=5.5cm]{figs/26_an2.png}
}

\memoanswer{
设粒子的入射速度为 $v$，第一次射出磁场的点为 $N_{0}'$，与板碰撞后再次进入磁场的位置为 $N_{1}$。粒子在磁场中运动的轨道半径为 $R$，有 $R=\frac{mv}{qB}$ …(1)；粒子速率不变，每次进入磁场与射出磁场位置间距离 $x_{1}$ 保持不变有 $x_{1}=N_{0}'N_{0}=2R\sin\theta$ …(2)；粒子射出磁场与下一次进入磁场位置间的距离 $x_{2}$ 始终不变，与 $N_{0}'N_{1}$ 相等，由图可以看出 $x_{2}=a$ …(3)。

设粒子最终离开磁场时，与挡板相碰 $n$ 次（$n=0$、1、2、3…）。若粒子能回到 P 点，由对称性，出射点的 $x$ 坐标应为 $-a$，即 $(n+1)x_{1}-nx_{2}=2a$ …(4)。由 (3)(4) 两式得 $x_{1}=\frac{n+2}{n+1}a$ …(5)。

若粒子与挡板发生碰撞，有 $x_{1}-x_{2}>\frac{a}{4}$ …(6)，联立 (3)(4)(6) 得 $n<3$ …(7)。联立 (1)(2)(5) 得 $v=\frac{qB}{2m\sin\theta}\cdot\frac{n+2}{n+1}a$ …(8)，把 $\sin\theta=\frac{h}{\sqrt{a^{2}+h^{2}}}$ 代入 (8) 中得 $v=\frac{qBa\sqrt{a^{2}+h^{2}}}{2mh}\cdot\frac{n+2}{n+1}$ …(9)。当 $n=0$ 时，$v_{0}=\frac{qBa\sqrt{a^{2}+h^{2}}}{mh}$ …(10)；当 $n=1$ 时，$v_{1}=\frac{3qBa\sqrt{a^{2}+h^{2}}}{4mh}$ …(11)；当 $n=2$ 时，$v_{2}=\frac{2qBa\sqrt{a^{2}+h^{2}}}{3mh}$ …(12)。

\onepicture[w=5.5cm]{figs/26_an1.png}
}

\end{enumerate}

\end{document}
